% Migrated from the Word baseline; do not duplicate content in other section files.
\section*{引言}
\label{sec:introduction}

磁驱柔性导管在微重力环境下同时受到磁场驱动、背景流体作用、血管壁接触及边界操控，其动力学呈现大转动、强几何非线性、弯曲---扭转耦合和开放端口并存等特征。若连续模型、有限维约化、数据表示与长期稳定性分析采用彼此不一致的变量或能量口径，模型误差、学习误差与数值离散误差将难以在同一理论框架内比较。因此，需要从连续力学、结构保持离散和可验证机器学习三个层面建立一致的数学接口。

几何精确梁杆理论为三维大变形柔性结构提供了统一的运动学与变分基础，磁力、流体载荷和接触模型则分别描述导管所处的多物理环境\cite{Simo1985FiniteStrainBeam,Antman2005Elasticity,Morison1950,COMSOLBeamInterface2025,COMSOLGeometricNonlinearity2025,COMSOLSectionStiffness2025,COMSOLWeakFormPDE2025,COMSOLWeakContribution2025,Sola2018MicroLie,Jackson1998Electrodynamics,Abbott2020MagneticRobotics,Newman1977MarineHydrodynamics,Faltinsen1990SeaLoads,Johnson1985ContactMechanics,HuntCrossley1975Restitution,BrownMcPhee2016Friction,MarsdenRatiu1999MechanicsSymmetry,HairerLubichWanner2006GNI}。然而，将平移与截面姿态联合保留、对保守与非保守载荷进行排他能量分账，并进一步获得适用于 Hamiltonian 学习和辛分析的有限维正则系统，仍需对弱形式、Galerkin 离散、约束约化和 Legendre 变换进行一致构造。

Hamiltonian 神经网络通过实值能量函数的辛梯度引入守恒归纳偏置，耗散与端口网络进一步扩展了受迫和耗散系统的结构化辨识；复值网络则可为幅值---相位耦合提供结构化表示\cite{Greydanus2019HNN,Zhong2020Dissipative,Desai2021PHNN,Chen2022KAMSLT,Trabelsi2018DeepComplex,Hirose2012CVNN}。但复值内部表示并不意味着物理 Hamiltonian 可以取一般复值，也不能由单一复标量无条件生成耗散和外部驱动。与此同时，多分支网络的结构分解不自动保证可辨识性，需要结合已知物理主体、端口功率监督、配对工况与消融实验进行验证。

KAM 理论与辛几何算法为近可积 Hamiltonian 系统及其数值映射的准周期结构提供了严格分析工具\cite{FengQin2010,Shang1999KAMAlgorithms,Shang2000KAMMappings,Haro2016Parameterization,Arnold1989ClassicalMechanics,Coquinot2020Metriplectic,Moser1962InvariantCurves,Kuksin1993InfiniteDimensionalKAM,Poschel2001KAMLecture}。经典连续 KAM、辛映射 KAM、二维扭转映射和无限维 KAM 分别对应不同的数学对象；训练损失或轨迹残差也不能直接替代高阶或解析函数范数中的小扰动条件。因此，机器学习误差必须与物理摄动、坐标变换、Galerkin 截断和时间离散误差在相容范数下统一计量。

本文的主要贡献概括如下：

（1）建立保留中心线平移、截面姿态、独立扭转和转动惯量的多物理几何精确导管模型，并以能量---功率恒等式统一磁、流体、接触与边界端口。

（2）通过平移---转动联合 Galerkin 离散、Kirchhoff 约束局部约化和一致 Legendre 变换，构造与 COMSOL 数据接口相容的有限维正则 Hamiltonian 系统。

（3）构造复值表示增强、实值 Hamiltonian 读出的结构化神经网络，并将保守修正、被动耗散、外部输入和一般残差分支化表示。

（4）建立统计学习误差、高阶正则性证书、辛离散误差与 KAM 条件之间的接口，并给出连续流、辛映射和后验不变环面核验的分支化分析框架。

本文的数值研究包括低维 P0 架构比较、三轨迹 P1/P1.5/P1.6 结构消融、COMSOL--Galerkin bridge 以及 canonical 数据合同诊断。现有有限证据至多支持单一 \(1\,\mathrm{s}\) 诊断哨兵作为后续 A1 资格实验的候选输入；它不认证跨 mode、幅值或工况的数据生成资格。计划的 96 条轨迹数据集、完整规模 P2 训练、开放端口辨识、高保真模型比较和数值 KAM 分析尚未完成，因此本文不据现有 pilot 结果推断正式的模型优势、不变环面或工程安全域。

全文其余部分安排如下：\Cref{sec:rod-model,sec:galerkin-discretization}建立连续多物理模型及联合 Galerkin 离散；\Cref{sec:hamiltonian-structure}构造有限维 Hamiltonian 结构与正则接口；\Cref{sec:chnn-model}给出复值表示增强的结构化 HNN 及统计误差分析；\Cref{sec:symplectic-kam}讨论辛离散、Birkhoff 结构和不变环面核验；\Cref{sec:numerical-results}用于报告数值实验，\Cref{sec:discussion}讨论适用范围与局限性，\Cref{sec:conclusion}总结全文。
